% =====================================================================
%  general_workflow.tex  --  include with % =====================================================================
%  general_workflow.tex  --  include with % =====================================================================
%  general_workflow.tex  --  include with % =====================================================================
%  general_workflow.tex  --  include with \input{general_workflow}
%  Requires in the MAIN preamble:
%     \usepackage{tikz}
%     \usetikzlibrary{arrows.meta,positioning,shapes.geometric,calc}
%     \usepackage{xcolor}   % loaded by tikz anyway
%  No \documentclass / \begin{document} here on purpose.
% =====================================================================

% ------------------------- colour palette ----------------------------
\definecolor{cInput}{HTML}{7B4B94}   % purple - input
\definecolor{cProc} {HTML}{2C6E9B}   % blue   - sequence / protein steps
\definecolor{cRNA}  {HTML}{1B7F5F}   % green  - RNA structure steps
\definecolor{cDec}  {HTML}{C1651A}   % orange - decision
\definecolor{cMan}  {HTML}{B3282B}   % red    - manual curation
\definecolor{cOut}  {HTML}{3F4A57}   % slate  - output
\definecolor{cSeed} {HTML}{8A8F98}   % grey   - reference data

% ------------------------------ styles -------------------------------
\tikzset{
  base/.style={draw, rounded corners=3pt, line width=0.9pt, align=center,
               inner sep=6pt, minimum height=13mm, text width=72mm,
               font=\small},
  input/.style ={base, draw=cInput, fill=cInput!8},
  proc/.style  ={base, draw=cProc,  fill=cProc!7},
  rna/.style   ={base, draw=cRNA,   fill=cRNA!7},
  manual/.style={base, draw=cMan,   fill=cMan!6, text width=60mm},
  output/.style={base, draw=cOut,   fill=cOut!10, line width=1.3pt},
  seed/.style  ={base, draw=cSeed,  fill=cSeed!12, text width=33mm,
                 minimum height=10mm, font=\footnotesize},
  decision/.style={draw=cDec, fill=cDec!8, line width=0.9pt, diamond,
                   aspect=2.2, align=center, inner sep=1pt,
                   text width=36mm, font=\small},
  flow/.style={-{Stealth[length=3mm,width=2mm]}, line width=0.9pt,
               draw=black!70},
  dashflow/.style={flow, dashed, dash pattern=on 3pt off 2pt, draw=cSeed},
  loop/.style={-{Stealth[length=3mm,width=2mm]}, line width=0.9pt,
               draw=cMan!85!black},
  elab/.style={font=\footnotesize\bfseries, inner sep=2pt, fill=white,
               text=black!75},
  note/.style={font=\scriptsize\itshape, text=black!60, align=left},
}

% \wfstep{number}{action}{tools}
\providecommand{\wfstep}[3]{\textbf{#1.~#2}\\[1.5pt]{\footnotesize\itshape #3}}

\begin{tikzpicture}[node distance=8mm]

% --------------------------- main chain ------------------------------
\node[input] (in)
  {\textbf{Input:} assembled contigs / genomes\\[1pt]
   {\footnotesize nucleotide FASTA --- isolates or metagenomes}};

\node[proc, below=of in] (orf)
  {\wfstep{1}{Gene calling / ORF prediction}
        {Prodigal $\rightarrow$ \texttt{.gff}, \texttt{.faa}, \texttt{.fna}}};

\node[proc, below=of orf] (rt)
  {\wfstep{2}{Identify the reverse transcriptase}
        {HMMER (\texttt{hmmsearch}) vs.\ retron-RT profiles}};

\node[proc, below=of rt] (win)
  {\wfstep{3}{Extract the genomic window up- and downstream of the RT}
        {intergenic regions $+$ flanking ORFs, $\pm$5--10\,kb}};

\node[rna, below=of win] (cm)
  {\wfstep{4}{Build covariance models from validated \textit{msr--msd} sequences}
        {Infernal \texttt{cmbuild} / \texttt{cmcalibrate}}};

\node[rna, below=of cm] (motif)
  {\wfstep{5}{Covariance-model-based ncRNA motif search in the non-coding region}
        {CMfinder $\cdot$ Infernal \texttt{cmsearch}}};

\node[decision, below=12mm of motif] (dec)
  {Structured\\ \textit{msr--msd}\\ recovered?};

\node[rna, below=12mm of dec] (align)
  {\wfstep{6}{Collect ncRNA candidate homologs and align}
        {MAFFT (\texttt{--Q-INS-i})}};

\node[rna, below=of align] (rscape)
  {\wfstep{7}{Validate covariation, control for phylogeny, confirm the functional structure}
        {R-scape}};

\node[proc, below=of rscape] (acc)
  {\wfstep{8}{Annotate the accessory / effector genes of the confirmed \textit{msr--msd}--RT unit}
        {InterPro $\cdot$ BLAST $\cdot$ HHpred / HHblits}};

\node[output, below=of acc] (out)
  {\textbf{Fully annotated retron operon}\\[2pt]
   {\footnotesize \textit{msr--msd} ncRNA \; $+$ \; RT \; $+$ \; accessory / effector gene(s)}};

% ------------------ reference data feeding the pipeline ---------------
\node[seed, left=15mm of rt] (seedRT)
  {\textbf{Reference set}\\ curated retron RT\\ protein sequences};
\node[seed, below=16mm of seedRT] (seedRNA)
  {\textbf{Reference set}\\ validated \textit{msr--msd}\\ transcripts $+$ relatives};

% ----------------------- manual curation branch -----------------------
\node[manual, right=18mm of dec, yshift=-4mm] (man)
  {\textbf{Manual curation}\\[2pt]
   {\footnotesize inspect the intergenic region for \textit{msr--msd}-like folds
   \emph{(RNAfold)}; build an RNA MSA of close relatives
   \emph{(MAFFT-Q-INS-i)}; trim the \textbf{a1}/\textbf{a2} arms
   and re-extract the flanks}};

% ------------------------------ arrows --------------------------------
\draw[flow] (in)     -- (orf);
\draw[flow] (orf)    -- (rt);
\draw[flow] (rt)     -- (win);
\draw[flow] (win)    -- (cm);
\draw[flow] (cm)     -- (motif);
\draw[flow] (motif)  -- (dec);
\draw[flow] (dec)    -- node[elab,right]{Yes} (align);
\draw[flow] (align)  -- (rscape);
\draw[flow] (rscape) -- (acc);
\draw[flow] (acc)    -- (out);

\draw[dashflow] (seedRT.east)  -- (rt.west);
\draw[dashflow] (seedRNA.east) -| ($(cm.west)+(-7mm,0)$) -- (cm.west);

% no-branch into manual curation
\draw[flow] (dec.east) -- node[elab,above,pos=0.4]{No} ++(9mm,0) |- (man.west);
% manual curation feeds the alignment step ...
\draw[loop] (man.south) |- (align.east);
% ... and iterates back to model building
\draw[loop] (man.north) -- ++(0,11mm) -| ($(cm.east)+(14mm,0)$)
            -- node[elab,pos=0.9,above]{iterate} (cm.east);

% ------------------------------- notes ---------------------------------
\node[note, anchor=north east, text width=45mm] at ($(cm.west)+(-14mm,2mm)$)
  {\textbf{Note:} sequence diversity usually\\ requires several independent CMs\\
   (one per retron clade / ncRNA family).};

\node[note, anchor=north, text width=105mm] at ($(out.south)+(0,-3mm)$)
  {\textbf{Note:} steps 4--7 are iterative --- trimming of the \textbf{a1}/\textbf{a2}
   regions and repeating the alignment typically improves the consensus
   structure before covariation is tested.};

\end{tikzpicture}
%  Requires in the MAIN preamble:
%     \usepackage{tikz}
%     \usetikzlibrary{arrows.meta,positioning,shapes.geometric,calc}
%     \usepackage{xcolor}   % loaded by tikz anyway
%  No \documentclass / \begin{document} here on purpose.
% =====================================================================

% ------------------------- colour palette ----------------------------
\definecolor{cInput}{HTML}{7B4B94}   % purple - input
\definecolor{cProc} {HTML}{2C6E9B}   % blue   - sequence / protein steps
\definecolor{cRNA}  {HTML}{1B7F5F}   % green  - RNA structure steps
\definecolor{cDec}  {HTML}{C1651A}   % orange - decision
\definecolor{cMan}  {HTML}{B3282B}   % red    - manual curation
\definecolor{cOut}  {HTML}{3F4A57}   % slate  - output
\definecolor{cSeed} {HTML}{8A8F98}   % grey   - reference data

% ------------------------------ styles -------------------------------
\tikzset{
  base/.style={draw, rounded corners=3pt, line width=0.9pt, align=center,
               inner sep=6pt, minimum height=13mm, text width=72mm,
               font=\small},
  input/.style ={base, draw=cInput, fill=cInput!8},
  proc/.style  ={base, draw=cProc,  fill=cProc!7},
  rna/.style   ={base, draw=cRNA,   fill=cRNA!7},
  manual/.style={base, draw=cMan,   fill=cMan!6, text width=60mm},
  output/.style={base, draw=cOut,   fill=cOut!10, line width=1.3pt},
  seed/.style  ={base, draw=cSeed,  fill=cSeed!12, text width=33mm,
                 minimum height=10mm, font=\footnotesize},
  decision/.style={draw=cDec, fill=cDec!8, line width=0.9pt, diamond,
                   aspect=2.2, align=center, inner sep=1pt,
                   text width=36mm, font=\small},
  flow/.style={-{Stealth[length=3mm,width=2mm]}, line width=0.9pt,
               draw=black!70},
  dashflow/.style={flow, dashed, dash pattern=on 3pt off 2pt, draw=cSeed},
  loop/.style={-{Stealth[length=3mm,width=2mm]}, line width=0.9pt,
               draw=cMan!85!black},
  elab/.style={font=\footnotesize\bfseries, inner sep=2pt, fill=white,
               text=black!75},
  note/.style={font=\scriptsize\itshape, text=black!60, align=left},
}

% \wfstep{number}{action}{tools}
\providecommand{\wfstep}[3]{\textbf{#1.~#2}\\[1.5pt]{\footnotesize\itshape #3}}

\begin{tikzpicture}[node distance=8mm]

% --------------------------- main chain ------------------------------
\node[input] (in)
  {\textbf{Input:} assembled contigs / genomes\\[1pt]
   {\footnotesize nucleotide FASTA --- isolates or metagenomes}};

\node[proc, below=of in] (orf)
  {\wfstep{1}{Gene calling / ORF prediction}
        {Prodigal $\rightarrow$ \texttt{.gff}, \texttt{.faa}, \texttt{.fna}}};

\node[proc, below=of orf] (rt)
  {\wfstep{2}{Identify the reverse transcriptase}
        {HMMER (\texttt{hmmsearch}) vs.\ retron-RT profiles}};

\node[proc, below=of rt] (win)
  {\wfstep{3}{Extract the genomic window up- and downstream of the RT}
        {intergenic regions $+$ flanking ORFs, $\pm$5--10\,kb}};

\node[rna, below=of win] (cm)
  {\wfstep{4}{Build covariance models from validated \textit{msr--msd} sequences}
        {Infernal \texttt{cmbuild} / \texttt{cmcalibrate}}};

\node[rna, below=of cm] (motif)
  {\wfstep{5}{Covariance-model-based ncRNA motif search in the non-coding region}
        {CMfinder $\cdot$ Infernal \texttt{cmsearch}}};

\node[decision, below=12mm of motif] (dec)
  {Structured\\ \textit{msr--msd}\\ recovered?};

\node[rna, below=12mm of dec] (align)
  {\wfstep{6}{Collect ncRNA candidate homologs and align}
        {MAFFT (\texttt{--Q-INS-i})}};

\node[rna, below=of align] (rscape)
  {\wfstep{7}{Validate covariation, control for phylogeny, confirm the functional structure}
        {R-scape}};

\node[proc, below=of rscape] (acc)
  {\wfstep{8}{Annotate the accessory / effector genes of the confirmed \textit{msr--msd}--RT unit}
        {InterPro $\cdot$ BLAST $\cdot$ HHpred / HHblits}};

\node[output, below=of acc] (out)
  {\textbf{Fully annotated retron operon}\\[2pt]
   {\footnotesize \textit{msr--msd} ncRNA \; $+$ \; RT \; $+$ \; accessory / effector gene(s)}};

% ------------------ reference data feeding the pipeline ---------------
\node[seed, left=15mm of rt] (seedRT)
  {\textbf{Reference set}\\ curated retron RT\\ protein sequences};
\node[seed, below=16mm of seedRT] (seedRNA)
  {\textbf{Reference set}\\ validated \textit{msr--msd}\\ transcripts $+$ relatives};

% ----------------------- manual curation branch -----------------------
\node[manual, right=18mm of dec, yshift=-4mm] (man)
  {\textbf{Manual curation}\\[2pt]
   {\footnotesize inspect the intergenic region for \textit{msr--msd}-like folds
   \emph{(RNAfold)}; build an RNA MSA of close relatives
   \emph{(MAFFT-Q-INS-i)}; trim the \textbf{a1}/\textbf{a2} arms
   and re-extract the flanks}};

% ------------------------------ arrows --------------------------------
\draw[flow] (in)     -- (orf);
\draw[flow] (orf)    -- (rt);
\draw[flow] (rt)     -- (win);
\draw[flow] (win)    -- (cm);
\draw[flow] (cm)     -- (motif);
\draw[flow] (motif)  -- (dec);
\draw[flow] (dec)    -- node[elab,right]{Yes} (align);
\draw[flow] (align)  -- (rscape);
\draw[flow] (rscape) -- (acc);
\draw[flow] (acc)    -- (out);

\draw[dashflow] (seedRT.east)  -- (rt.west);
\draw[dashflow] (seedRNA.east) -| ($(cm.west)+(-7mm,0)$) -- (cm.west);

% no-branch into manual curation
\draw[flow] (dec.east) -- node[elab,above,pos=0.4]{No} ++(9mm,0) |- (man.west);
% manual curation feeds the alignment step ...
\draw[loop] (man.south) |- (align.east);
% ... and iterates back to model building
\draw[loop] (man.north) -- ++(0,11mm) -| ($(cm.east)+(14mm,0)$)
            -- node[elab,pos=0.9,above]{iterate} (cm.east);

% ------------------------------- notes ---------------------------------
\node[note, anchor=north east, text width=45mm] at ($(cm.west)+(-14mm,2mm)$)
  {\textbf{Note:} sequence diversity usually\\ requires several independent CMs\\
   (one per retron clade / ncRNA family).};

\node[note, anchor=north, text width=105mm] at ($(out.south)+(0,-3mm)$)
  {\textbf{Note:} steps 4--7 are iterative --- trimming of the \textbf{a1}/\textbf{a2}
   regions and repeating the alignment typically improves the consensus
   structure before covariation is tested.};

\end{tikzpicture}
%  Requires in the MAIN preamble:
%     \usepackage{tikz}
%     \usetikzlibrary{arrows.meta,positioning,shapes.geometric,calc}
%     \usepackage{xcolor}   % loaded by tikz anyway
%  No \documentclass / \begin{document} here on purpose.
% =====================================================================

% ------------------------- colour palette ----------------------------
\definecolor{cInput}{HTML}{7B4B94}   % purple - input
\definecolor{cProc} {HTML}{2C6E9B}   % blue   - sequence / protein steps
\definecolor{cRNA}  {HTML}{1B7F5F}   % green  - RNA structure steps
\definecolor{cDec}  {HTML}{C1651A}   % orange - decision
\definecolor{cMan}  {HTML}{B3282B}   % red    - manual curation
\definecolor{cOut}  {HTML}{3F4A57}   % slate  - output
\definecolor{cSeed} {HTML}{8A8F98}   % grey   - reference data

% ------------------------------ styles -------------------------------
\tikzset{
  base/.style={draw, rounded corners=3pt, line width=0.9pt, align=center,
               inner sep=6pt, minimum height=13mm, text width=72mm,
               font=\small},
  input/.style ={base, draw=cInput, fill=cInput!8},
  proc/.style  ={base, draw=cProc,  fill=cProc!7},
  rna/.style   ={base, draw=cRNA,   fill=cRNA!7},
  manual/.style={base, draw=cMan,   fill=cMan!6, text width=60mm},
  output/.style={base, draw=cOut,   fill=cOut!10, line width=1.3pt},
  seed/.style  ={base, draw=cSeed,  fill=cSeed!12, text width=33mm,
                 minimum height=10mm, font=\footnotesize},
  decision/.style={draw=cDec, fill=cDec!8, line width=0.9pt, diamond,
                   aspect=2.2, align=center, inner sep=1pt,
                   text width=36mm, font=\small},
  flow/.style={-{Stealth[length=3mm,width=2mm]}, line width=0.9pt,
               draw=black!70},
  dashflow/.style={flow, dashed, dash pattern=on 3pt off 2pt, draw=cSeed},
  loop/.style={-{Stealth[length=3mm,width=2mm]}, line width=0.9pt,
               draw=cMan!85!black},
  elab/.style={font=\footnotesize\bfseries, inner sep=2pt, fill=white,
               text=black!75},
  note/.style={font=\scriptsize\itshape, text=black!60, align=left},
}

% \wfstep{number}{action}{tools}
\providecommand{\wfstep}[3]{\textbf{#1.~#2}\\[1.5pt]{\footnotesize\itshape #3}}

\begin{tikzpicture}[node distance=8mm]

% --------------------------- main chain ------------------------------
\node[input] (in)
  {\textbf{Input:} assembled contigs / genomes\\[1pt]
   {\footnotesize nucleotide FASTA --- isolates or metagenomes}};

\node[proc, below=of in] (orf)
  {\wfstep{1}{Gene calling / ORF prediction}
        {Prodigal $\rightarrow$ \texttt{.gff}, \texttt{.faa}, \texttt{.fna}}};

\node[proc, below=of orf] (rt)
  {\wfstep{2}{Identify the reverse transcriptase}
        {HMMER (\texttt{hmmsearch}) vs.\ retron-RT profiles}};

\node[proc, below=of rt] (win)
  {\wfstep{3}{Extract the genomic window up- and downstream of the RT}
        {intergenic regions $+$ flanking ORFs, $\pm$5--10\,kb}};

\node[rna, below=of win] (cm)
  {\wfstep{4}{Build covariance models from validated \textit{msr--msd} sequences}
        {Infernal \texttt{cmbuild} / \texttt{cmcalibrate}}};

\node[rna, below=of cm] (motif)
  {\wfstep{5}{Covariance-model-based ncRNA motif search in the non-coding region}
        {CMfinder $\cdot$ Infernal \texttt{cmsearch}}};

\node[decision, below=12mm of motif] (dec)
  {Structured\\ \textit{msr--msd}\\ recovered?};

\node[rna, below=12mm of dec] (align)
  {\wfstep{6}{Collect ncRNA candidate homologs and align}
        {MAFFT (\texttt{--Q-INS-i})}};

\node[rna, below=of align] (rscape)
  {\wfstep{7}{Validate covariation, control for phylogeny, confirm the functional structure}
        {R-scape}};

\node[proc, below=of rscape] (acc)
  {\wfstep{8}{Annotate the accessory / effector genes of the confirmed \textit{msr--msd}--RT unit}
        {InterPro $\cdot$ BLAST $\cdot$ HHpred / HHblits}};

\node[output, below=of acc] (out)
  {\textbf{Fully annotated retron operon}\\[2pt]
   {\footnotesize \textit{msr--msd} ncRNA \; $+$ \; RT \; $+$ \; accessory / effector gene(s)}};

% ------------------ reference data feeding the pipeline ---------------
\node[seed, left=15mm of rt] (seedRT)
  {\textbf{Reference set}\\ curated retron RT\\ protein sequences};
\node[seed, below=16mm of seedRT] (seedRNA)
  {\textbf{Reference set}\\ validated \textit{msr--msd}\\ transcripts $+$ relatives};

% ----------------------- manual curation branch -----------------------
\node[manual, right=18mm of dec, yshift=-4mm] (man)
  {\textbf{Manual curation}\\[2pt]
   {\footnotesize inspect the intergenic region for \textit{msr--msd}-like folds
   \emph{(RNAfold)}; build an RNA MSA of close relatives
   \emph{(MAFFT-Q-INS-i)}; trim the \textbf{a1}/\textbf{a2} arms
   and re-extract the flanks}};

% ------------------------------ arrows --------------------------------
\draw[flow] (in)     -- (orf);
\draw[flow] (orf)    -- (rt);
\draw[flow] (rt)     -- (win);
\draw[flow] (win)    -- (cm);
\draw[flow] (cm)     -- (motif);
\draw[flow] (motif)  -- (dec);
\draw[flow] (dec)    -- node[elab,right]{Yes} (align);
\draw[flow] (align)  -- (rscape);
\draw[flow] (rscape) -- (acc);
\draw[flow] (acc)    -- (out);

\draw[dashflow] (seedRT.east)  -- (rt.west);
\draw[dashflow] (seedRNA.east) -| ($(cm.west)+(-7mm,0)$) -- (cm.west);

% no-branch into manual curation
\draw[flow] (dec.east) -- node[elab,above,pos=0.4]{No} ++(9mm,0) |- (man.west);
% manual curation feeds the alignment step ...
\draw[loop] (man.south) |- (align.east);
% ... and iterates back to model building
\draw[loop] (man.north) -- ++(0,11mm) -| ($(cm.east)+(14mm,0)$)
            -- node[elab,pos=0.9,above]{iterate} (cm.east);

% ------------------------------- notes ---------------------------------
\node[note, anchor=north east, text width=45mm] at ($(cm.west)+(-14mm,2mm)$)
  {\textbf{Note:} sequence diversity usually\\ requires several independent CMs\\
   (one per retron clade / ncRNA family).};

\node[note, anchor=north, text width=105mm] at ($(out.south)+(0,-3mm)$)
  {\textbf{Note:} steps 4--7 are iterative --- trimming of the \textbf{a1}/\textbf{a2}
   regions and repeating the alignment typically improves the consensus
   structure before covariation is tested.};

\end{tikzpicture}
%  Requires in the MAIN preamble:
%     \usepackage{tikz}
%     \usetikzlibrary{arrows.meta,positioning,shapes.geometric,calc}
%     \usepackage{xcolor}   % loaded by tikz anyway
%  No \documentclass / \begin{document} here on purpose.
% =====================================================================

% ------------------------- colour palette ----------------------------
\definecolor{cInput}{HTML}{7B4B94}   % purple - input
\definecolor{cProc} {HTML}{2C6E9B}   % blue   - sequence / protein steps
\definecolor{cRNA}  {HTML}{1B7F5F}   % green  - RNA structure steps
\definecolor{cDec}  {HTML}{C1651A}   % orange - decision
\definecolor{cMan}  {HTML}{B3282B}   % red    - manual curation
\definecolor{cOut}  {HTML}{3F4A57}   % slate  - output
\definecolor{cSeed} {HTML}{8A8F98}   % grey   - reference data

% ------------------------------ styles -------------------------------
\tikzset{
  base/.style={draw, rounded corners=3pt, line width=0.9pt, align=center,
               inner sep=6pt, minimum height=13mm, text width=72mm,
               font=\small},
  input/.style ={base, draw=cInput, fill=cInput!8},
  proc/.style  ={base, draw=cProc,  fill=cProc!7},
  rna/.style   ={base, draw=cRNA,   fill=cRNA!7},
  manual/.style={base, draw=cMan,   fill=cMan!6, text width=60mm},
  output/.style={base, draw=cOut,   fill=cOut!10, line width=1.3pt},
  seed/.style  ={base, draw=cSeed,  fill=cSeed!12, text width=33mm,
                 minimum height=10mm, font=\footnotesize},
  decision/.style={draw=cDec, fill=cDec!8, line width=0.9pt, diamond,
                   aspect=2.2, align=center, inner sep=1pt,
                   text width=36mm, font=\small},
  flow/.style={-{Stealth[length=3mm,width=2mm]}, line width=0.9pt,
               draw=black!70},
  dashflow/.style={flow, dashed, dash pattern=on 3pt off 2pt, draw=cSeed},
  loop/.style={-{Stealth[length=3mm,width=2mm]}, line width=0.9pt,
               draw=cMan!85!black},
  elab/.style={font=\footnotesize\bfseries, inner sep=2pt, fill=white,
               text=black!75},
  note/.style={font=\scriptsize\itshape, text=black!60, align=left},
}

% \wfstep{number}{action}{tools}
\providecommand{\wfstep}[3]{\textbf{#1.~#2}\\[1.5pt]{\footnotesize\itshape #3}}

\begin{tikzpicture}[node distance=8mm]

% --------------------------- main chain ------------------------------
\node[input] (in)
  {\textbf{Input:} assembled contigs / genomes\\[1pt]
   {\footnotesize nucleotide FASTA --- isolates or metagenomes}};

\node[proc, below=of in] (orf)
  {\wfstep{1}{Gene calling / ORF prediction}
        {Prodigal $\rightarrow$ \texttt{.gff}, \texttt{.faa}, \texttt{.fna}}};

\node[proc, below=of orf] (rt)
  {\wfstep{2}{Identify the reverse transcriptase}
        {HMMER (\texttt{hmmsearch}) vs.\ retron-RT profiles}};

\node[proc, below=of rt] (win)
  {\wfstep{3}{Extract the genomic window up- and downstream of the RT}
        {intergenic regions $+$ flanking ORFs, $\pm$5--10\,kb}};

\node[rna, below=of win] (cm)
  {\wfstep{4}{Build covariance models from validated \textit{msr--msd} sequences}
        {Infernal \texttt{cmbuild} / \texttt{cmcalibrate}}};

\node[rna, below=of cm] (motif)
  {\wfstep{5}{Covariance-model-based ncRNA motif search in the non-coding region}
        {CMfinder $\cdot$ Infernal \texttt{cmsearch}}};

\node[decision, below=12mm of motif] (dec)
  {Structured\\ \textit{msr--msd}\\ recovered?};

\node[rna, below=12mm of dec] (align)
  {\wfstep{6}{Collect ncRNA candidate homologs and align}
        {MAFFT (\texttt{--Q-INS-i})}};

\node[rna, below=of align] (rscape)
  {\wfstep{7}{Validate covariation, control for phylogeny, confirm the functional structure}
        {R-scape}};

\node[proc, below=of rscape] (acc)
  {\wfstep{8}{Annotate the accessory / effector genes of the confirmed \textit{msr--msd}--RT unit}
        {InterPro $\cdot$ BLAST $\cdot$ HHpred / HHblits}};

\node[output, below=of acc] (out)
  {\textbf{Fully annotated retron operon}\\[2pt]
   {\footnotesize \textit{msr--msd} ncRNA \; $+$ \; RT \; $+$ \; accessory / effector gene(s)}};

% ------------------ reference data feeding the pipeline ---------------
\node[seed, left=15mm of rt] (seedRT)
  {\textbf{Reference set}\\ curated retron RT\\ protein sequences};
\node[seed, below=16mm of seedRT] (seedRNA)
  {\textbf{Reference set}\\ validated \textit{msr--msd}\\ transcripts $+$ relatives};

% ----------------------- manual curation branch -----------------------
\node[manual, right=18mm of dec, yshift=-4mm] (man)
  {\textbf{Manual curation}\\[2pt]
   {\footnotesize inspect the intergenic region for \textit{msr--msd}-like folds
   \emph{(RNAfold)}; build an RNA MSA of close relatives
   \emph{(MAFFT-Q-INS-i)}; trim the \textbf{a1}/\textbf{a2} arms
   and re-extract the flanks}};

% ------------------------------ arrows --------------------------------
\draw[flow] (in)     -- (orf);
\draw[flow] (orf)    -- (rt);
\draw[flow] (rt)     -- (win);
\draw[flow] (win)    -- (cm);
\draw[flow] (cm)     -- (motif);
\draw[flow] (motif)  -- (dec);
\draw[flow] (dec)    -- node[elab,right]{Yes} (align);
\draw[flow] (align)  -- (rscape);
\draw[flow] (rscape) -- (acc);
\draw[flow] (acc)    -- (out);

\draw[dashflow] (seedRT.east)  -- (rt.west);
\draw[dashflow] (seedRNA.east) -| ($(cm.west)+(-7mm,0)$) -- (cm.west);

% no-branch into manual curation
\draw[flow] (dec.east) -- node[elab,above,pos=0.4]{No} ++(9mm,0) |- (man.west);
% manual curation feeds the alignment step ...
\draw[loop] (man.south) |- (align.east);
% ... and iterates back to model building
\draw[loop] (man.north) -- ++(0,11mm) -| ($(cm.east)+(14mm,0)$)
            -- node[elab,pos=0.9,above]{iterate} (cm.east);

% ------------------------------- notes ---------------------------------
\node[note, anchor=north east, text width=45mm] at ($(cm.west)+(-14mm,2mm)$)
  {\textbf{Note:} sequence diversity usually\\ requires several independent CMs\\
   (one per retron clade / ncRNA family).};

\node[note, anchor=north, text width=105mm] at ($(out.south)+(0,-3mm)$)
  {\textbf{Note:} steps 4--7 are iterative --- trimming of the \textbf{a1}/\textbf{a2}
   regions and repeating the alignment typically improves the consensus
   structure before covariation is tested.};

\end{tikzpicture}